\chapter*{前言}
\addcontentsline{toc}{chapter}{前言}
\markboth{前言}{前言}

\section*{为什么写这本笔记}

过去十年，\kw{流量}这门生意被彻底重构成了一个\kw{算法 + 数据 + 供应链}的复合体：
搜索引擎把网页变成倒排索引与排序模型，推荐系统把「人找货」变成「货找人」，
计算广告把每一次曝光放进实时拍卖里定价，大语言模型（LLM）又把「检索—生成」的边界重新画了一遍。
与此同时，Temu、SHEIN、TikTok Shop 这类平台把中国供应链的产能，用极致的数据化方式直接怼到海外消费者面前。

这五条线索表面分散，底层其实是同一套东西：

\begin{keypoint}[一句话总纲]
\textbf{在约束条件下，用数据把「谁—在什么场景—看到什么—以什么价格—由谁履约」这条链路逐步逼近最优。}
推荐、搜索、广告解决的是\textbf{匹配}问题；SEO/GEO/投流解决的是\textbf{获取}问题；
跨境电商与供应链解决的是\textbf{交付与成本}问题。三者共享同一套指标体系、同一套实验方法、同一套特征与排序基础设施。
\end{keypoint}

这本笔记的目标不是「讲完所有知识」，而是搭一张\textbf{可以持续挂东西的骨架}：
每一章先给出问题定义与指标体系，再给出主流方法的技术演进脉络，然后落到工程实现、
常见陷阱与真实业务中的取舍，最后附上可追的论文与工具。

\section*{这本笔记适合谁}

\begin{itemize}
  \item 准备进入或已经在大厂做\textbf{推荐 / 搜索 / 广告}算法与工程的工程师；
  \item 负责\textbf{自然流量（SEO/GEO）与付费流量（Google/Meta/TikTok）}的增长负责人、独立站操盘手；
  \item 做\textbf{跨境电商}、需要理解平台规则、物流履约与供应链选品的卖家与运营；
  \item 想把这五块内容串成一条完整链路、形成\textbf{系统认知}的转行者与在校生。
\end{itemize}

\section*{使用约定}

\begin{description}[leftmargin=5em, style=nextline]
  \item[术语] 首次出现时给出中英对照，例如 \term{协同过滤}{Collaborative Filtering}；完整对照表见\cref{app:glossary}。
  \item[公式] 尽量给出\textbf{业务含义}与\textbf{可实现的简化形式}，而不是只堆符号。
  \item[案例] 标注为「案例」的方框，多数来自公开资料、论文实验或行业通识，用于说明取舍逻辑，不代表任何公司的内部数据。
  \item[数据] 文中提到的具体费率、CPM、转化率等均为\textbf{示范量级}，实际投放请以账户后台与最新官方文档为准。
  \item[代码] 代码块以 Python / SQL / 配置片段为主，仅用于说明思路，未做生产级健壮性处理。
\end{description}

\section*{免责声明}

广告平台政策、平台招商规则、税务与合规要求（VAT、EPR、DSA、GPSR、UFLPA 等）变动极快，
本笔记中的相关描述仅反映写作时的通行做法，\textbf{不构成法律、税务或投资建议}。
涉及合规决策，请以官方文件与专业顾问意见为准。

\vspace{1cm}
\begin{flushright}
\textit{—— 愿你既懂模型，也懂生意。}
\end{flushright}
